\section{Exact dimensions for Haar-typical multipliers}
\label{typical:section}

For a unit $\lambda\in\Z_3^\times$ put
\[
 C_\lambda=\Sigma\cap\lambda^{-1}\Sigma,
 \qquad H=1+3\Z_3,
 \qquad S_\lambda=C_\lambda\cap H,
 \qquad b=\log_3(4/3).
\]
Write $m$ for normalized additive Haar measure on $\Z_3$ and $m_H$
for its normalized restriction to $H$. Abram and Lagarias
\cite[Section~2.4]{AL2014} propose studying the dimension function
$\lambda\mapsto\dimH C_\lambda$ for general $3$-adic parameters,
including its typical behavior. The next theorem determines its
Haar-typical value on each of the two residue classes of units.
It does not determine all exceptional parameters or the Baire-typical
value.

\begin{theorem}[Exact Haar-typical dimension]\label{typical:main}
For $m_H$-almost every $\lambda$, simultaneously for
$K=C_\lambda$ and $K=S_\lambda$,
\begin{equation}\label{typical:four-dimensions}
 \dimH K=\dimP K=\underline{\dim}_{\mathrm B}K
       =\overline{\dim}_{\mathrm B}K=\log_3(4/3).
\end{equation}
For every $\lambda\in2+3\Z_3$ one has $C_\lambda=\{0\}$.
\end{theorem}

The proof computes an averaged second moment exactly, proves positivity
of the resulting affine density, and transfers affine block entropy
to the homogeneous intersection. The last transfer is separate from
the almost-everywhere assertion about the affine translation: one
cannot obtain it by setting that translation equal to zero.
Throughout this section logarithms without a subscript are natural.
For $z\in\Z_3$, let $[z]_n\in\{0,\ldots,3^n-1\}$ be its canonical
residue modulo $3^n$, and let $\Sigma_n$ be the $2^n$ canonical
residues with digits $0,1$.

\subsection{Affine fibers and their exact second moment}

For $\lambda\in H$ define
\begin{equation}\label{typical:affine-count}
 \begin{split}
 R_n(\lambda,\beta)
   &=\#\{x\in\Sigma_n:[\lambda x+\beta]_n\in\Sigma_n\},\\
 H_n(\lambda)&=3^{-n}\sum_{\beta\bmod3^n}\log R_n(\lambda,\beta).
 \end{split}
\end{equation}
For every $\beta\in\Z_3$,
\begin{equation}\label{typical:count-range}
 1\le R_n(\lambda,\beta)\le2^n,
 \qquad
 3^{-n}\sum_{\beta\bmod3^n}R_n(\lambda,\beta)=(4/3)^n.
\end{equation}
Indeed, after an admissible input prefix the next output digit is
$c+a$ modulo $3$, where $a\in\{0,1\}$. At least one of $c,c+1$
belongs to $\{0,1\}$ modulo $3$. This proves both nonemptiness and
extension of every admissible finite prefix to an infinite point.
For the average, each of the $4^n$ ordered pairs of input and output
words determines a unique translation modulo $3^n$.

Let $\mu$ be the probability measure on $\Sigma$ with independent
fair digits, and let $\nu_\lambda$ be the law of $Y-\lambda X$ for
independent $X,Y$ with law $\mu$. Its cylinder density is
\begin{equation}\label{typical:cylinder-density}
 D_{\lambda,n}(\beta)
 =3^n\nu_\lambda(\beta+3^n\Z_3)
 =(3/4)^nR_n(\lambda,\beta).
\end{equation}

\begin{lemma}[Exact averaged second moment]\label{typical:l2}
For every $n\ge1$,
\begin{equation}\label{typical:l2-identity}
 \int_H\|D_{\lambda,n}\|_{L^2(m)}^2\,dm_H(\lambda)
       =\frac32-\frac12\left(\frac34\right)^n.
\end{equation}
Consequently $\nu_\lambda$ has an $L^2(m)$ density for
$m_H$-almost every $\lambda$.
\end{lemma}

\begin{proof}
Take independent $X,X',Y,Y'$ with law $\mu$, and set
$A=X-X'$, $B=Y-Y'$. Counting coincident residue classes gives
\[
 \|D_{\lambda,n}\|_2^2
       =3^n\mathbb P(B\equiv\lambda A\pmod{3^n}).
\]
The first differing digit of $X,X'$ gives
\[
 \mathbb P(v_3(A)\ge n)=2^{-n},\qquad
 \mathbb P(v_3(A)=k)=2^{-k-1}\quad(k\ge0).
\]
Conditioned on $v_3(A)=k$, the first unit digit of $A/3^k$
is $+1$ or $-1$ with equal probabilities; the analogous digit of
$B/3^k$ is independent. If both valuations are at least $n$, the
congruence holds. If both equal $k<n$, their leading signs agree
with probability $1/2$. On this event the required residue of
$\lambda\in H$ modulo $3^{n-k}$ has measure $3^{-(n-k-1)}$.
Unequal valuations below $n$ cannot contribute. Hence the average is
\begin{align*}
 3^n\left[4^{-n}+
   \sum_{k=0}^{n-1}4^{-k-1}\frac12\,3^{-(n-k-1)}\right]
 &=\left(\frac34\right)^n+
   \frac38\sum_{k=0}^{n-1}\left(\frac34\right)^k\\
 &=\frac32-\frac12\left(\frac34\right)^n.
\end{align*}
For fixed $\lambda$, the functions $D_{\lambda,n}$ are a
nonnegative martingale for the increasing cylinder sigma-fields of
$\beta$ under $m$. Their squared $L^2$ norms increase. Monotone
convergence and \eqref{typical:l2-identity} yield
\[
 \int_H\sup_n\|D_{\lambda,n}\|_2^2\,dm_H(\lambda)=\frac32.
\]
For almost every $\lambda$ the martingale is therefore bounded in
$L^2$ and converges in $L^2$ and almost everywhere to some $D_\lambda$.
Its integrals on every cylinder agree with $\nu_\lambda$; since
cylinders generate the Borel sigma-field, $\nu_\lambda=D_\lambda m$.
In particular $\int D_\lambda\,dm=1$.
\end{proof}

\begin{lemma}[Positivity of an affine density]\label{typical:positive}
Fix $\lambda\in H$. If $\nu_\lambda\ll m$, its density is positive
$m$-almost everywhere.
\end{lemma}

\begin{proof}
Decomposing the first input and output digits gives
\[
 \nu_\lambda=\frac14
  \sum_{d\in(0,1,-\lambda,1-\lambda)}(z\mapsto d+3z)_*\nu_\lambda,
\]
where the four terms retain their multiplicities. A nonnegative
version $D$ of the density consequently satisfies, almost everywhere,
\begin{equation}\label{typical:density-relation}
 D(z)=\frac34\sum_{\substack{d\in(0,1,-\lambda,1-\lambda)\\
                            d\equiv z\ (3)}}D((z-d)/3).
\end{equation}
The three representatives $d_0=0,d_1=1,d_2=-\lambda$ lie in
different residue classes. Define $T(z)=(z-d_r)/3$ when
$z\equiv r\pmod3$. Each inverse branch scales Haar measure by $1/3$,
so $T$ preserves $m$. It is also ergodic. In fact,
\[
 (r_0,r_1,\ldots)\longmapsto\sum_{j\ge0}3^j d_{r_j}
\]
is an isometric bijection from the full ternary sequence space onto
$\Z_3$. Distinct representatives modulo $3$ preserve the first
differing position, and repeated application of $T$ gives the inverse
coding. The remainder $3^nT^n(z)$ tends to zero. Each sequence
cylinder maps onto a residue class of the same measure, so this
coding sends uniform Bernoulli measure to $m$ and conjugates $T$
to the ergodic one-sided shift.

For $Z=\{D=0\}$, nonnegativity in
\eqref{typical:density-relation} implies $Z\subseteq T^{-1}Z$
modulo null sets. The two sets have the same measure, hence agree
modulo null sets. Ergodicity gives $m(Z)\in\{0,1\}$, and
$\int D\,dm=1$ excludes $m(Z)=1$. All uses of almost-everywhere
identities are legitimate because the inverse branches identify Haar
measure with its normalized restrictions to residue classes.
\end{proof}

\begin{corollary}[Saturation of affine block entropy]
\label{typical:entropy-limit}
For $m_H$-almost every $\lambda$,
\begin{equation}\label{typical:entropy-limit-eq}
 \lim_{n\to\infty}\frac{H_n(\lambda)}n=\log(4/3).
\end{equation}
\end{corollary}

\begin{proof}
For almost every $\lambda$, Lemmas~\ref{typical:l2} and
\ref{typical:positive} give $0<D_\lambda(\beta)<\infty$ for
almost every $\beta$, with $D_{\lambda,n}(\beta)\to D_\lambda(\beta)$.
Thus $n^{-1}\log D_{\lambda,n}(\beta)\to0$. By
\eqref{typical:count-range} and \eqref{typical:cylinder-density},
\[
 -\log(4/3)\le\frac1n\log D_{\lambda,n}(\beta)
              \le\log(3/2).
\]
Bounded convergence now gives
\[
 \frac{H_n(\lambda)}n
 =\log(4/3)+\int\frac1n\log D_{\lambda,n}\,dm
 \longrightarrow\log(4/3).
\]
No integrability assumption on $\log D_\lambda$ is needed.
\end{proof}

\subsection{Transfer to the homogeneous fiber}

\begin{lemma}[Block measures on a prescribed homogeneous fiber]
\label{typical:block-measure}
For each fixed integer $q\ge1$ and $m_H$-almost every $\lambda$,
there is an exact-dimensional probability measure on $S_\lambda$
of dimension $H_q(\lambda)/(q\log3)$. In particular,
\begin{equation}\label{typical:block-lower}
 \dimH S_\lambda\ge \frac{H_q(\lambda)}{q\log3}.
\end{equation}
\end{lemma}

\begin{proof}
Partition $H$ into its finitely many residue classes modulo $3^q$.
In one such class fix $\ell=[\lambda]_q$. This also fixes
$H_q(\lambda)=H_q(\ell)$. Choose an admissible unit prefix
$P_q\in\Sigma_q$ deterministically from $\ell$: start with digit $1$
and take the least admissible digit thereafter. The extension property
proved after \eqref{typical:count-range} ensures this choice exists.

Suppose an admissible unit prefix $P_j$ of length $j=tq$ has been
chosen. An input block $z\in\Sigma_q$ produces the output block
\begin{equation}\label{typical:block-equation}
 [\ell z+\beta_j]_q,
 \qquad
 \beta_j=\left[\frac{\lambda P_j-[\lambda P_j]_j}{3^j}\right]_q.
\end{equation}
There are exactly $R_q(\ell,\beta_j)$ admissible blocks. Give them
equal probabilities and set $P_{j+q}=P_j+3^jz$. For each fixed
$\lambda$ these consistent cylinder probabilities define a
probability measure $\eta_{\lambda,q}$ supported on $S_\lambda$.
They depend on finitely many parameter digits at each finite stage,
so $\lambda\mapsto\eta_{\lambda,q}$ is a measurable probability kernel.

To analyze the measure, first choose $\lambda$ with conditional Haar
law in the class $\ell$, then implement the block choices using
independent uniform random variables $V_1,V_2,\ldots$. Before the
choice at $j=tq$, the past sigma-field is
\[
 \mathcal F_t=\sigma([\lambda]_{tq},V_1,\ldots,V_{t-1}).
\]
The prefix $P_j$ is measurable with respect to this past. The fresh
parameter block $L_j=[(\lambda-[\lambda]_j)/3^j]_q$ is uniform modulo
$3^q$ and independent of $\mathcal F_t$. Equation
\eqref{typical:block-equation} gives
\begin{equation}\label{typical:fresh-shift}
 \beta_j\equiv c_j+L_jP_j\pmod{3^q},
\end{equation}
where $c_j$ is determined by the past. Since $P_j\equiv1\pmod3$,
multiplication by $P_j$ permutes all residues modulo $3^q$.
Thus $\beta_j$ has the uniform conditional law for every past.
After using $V_t$, the next prefix is measurable with respect to
$\mathcal F_{t+1}$, which proves the induction.

Consequently $Y_t=\log R_q(\ell,\beta_{tq})$ are independent and
identically distributed in this joint experiment, with
$0\le Y_t\le q\log2$ and expectation $H_q(\ell)$. At the block
boundary $n=q(T+1)$ the mass of the actual cylinder is exactly
\[
 \eta_{\lambda,q}(x+3^n\Z_3)
       =\prod_{t=1}^T R_q(\ell,\beta_{tq})^{-1}.
\]
The strong law gives local dimension $H_q(\ell)/(q\log3)$ along
these boundaries, jointly almost surely. Conditional on the full
parameter $\lambda$, the $V_t$ still implement exactly
$\eta_{\lambda,q}$. Fubini's theorem therefore gives the same
almost-sure local dimension for almost every fixed $\lambda$.
The parameter marginal remains Haar throughout; no independence
of the shifts for a fixed $\lambda$ is asserted.

Between consecutive block boundaries a cylinder mass is bounded
by the two boundary masses. Their ratio is at most $2^q$, and the
depths differ by $q$. Dividing logarithms by $n\log3$ shows that
the local dimension exists on all scales. Exact dimensionality
implies \eqref{typical:block-lower}: for any smaller exponent,
restrict the measure to a positive-measure set on which the
corresponding cylinder upper bound holds beyond one common depth,
and apply the mass distribution principle. Taking the union of the
finitely many residue classes completes the proof.
\end{proof}

Intersect the full-measure parameter sets from
Lemma~\ref{typical:block-measure} for every $q\ge1$, and also the
full-measure set of Corollary~\ref{typical:entropy-limit}. On their
intersection,
\begin{equation}\label{typical:hausdorff-lower}
 \dimH S_\lambda\ge
 \sup_{q\ge1}\frac{H_q(\lambda)}{q\log3}\ge b.
\end{equation}
This is the required passage to translation zero. It uses fresh
parameter blocks rather than a substitution into an almost-everywhere
statement in the translation variable.

\subsection{The upper bound and the initial zero digits}

Let $Q_n(\lambda)$ be the number of admissible unit prefixes of
length $n$. Every such prefix extends, so this is the exact number
of depth-$n$ cylinders meeting $S_\lambda$. For each of the
$2^{n-1}$ possible unit input words $P$, the residue $\lambda P$
is uniform among the $3^{n-1}$ residues congruent to $1$ modulo $3$
when $\lambda$ has law $m_H$. Exactly $2^{n-1}$ of those residues
have digits $0,1$. Hence
\begin{equation}\label{typical:unit-first-moment}
 \mathbb E_H Q_n=(4/3)^{n-1}.
\end{equation}
For every $\varepsilon>0$, Markov's inequality gives
\[
 m_H\{Q_n>3^{n(b+\varepsilon)}\}
     \le\tfrac34\,3^{-n\varepsilon}.
\]
This is summable. Borel--Cantelli, followed by a countable
intersection over positive rational $\varepsilon$, shows that
almost surely $\overline{\dim}_{\mathrm B}S_\lambda\le b$.

To handle $C_\lambda$, count the possible position of the first digit
$1$. With $N_n(C_\lambda)$ denoting its number of depth-$n$ cylinders,
\begin{equation}\label{typical:full-prefix-count}
 N_n(C_\lambda)=1+\sum_{j=1}^n Q_j(\lambda).
\end{equation}
The first term is the all-zero prefix. A prefix with $v<n$ initial
zeros corresponds to a unit prefix of length $n-v$. For almost every
$\lambda$ and every $\varepsilon>0$, the finitely many exceptions
can be absorbed into a finite constant:
$Q_j(\lambda)\le K_{\lambda,\varepsilon}3^{j(b+\varepsilon)}$
for all $j$. A geometric sum in \eqref{typical:full-prefix-count}
therefore gives $\overline{\dim}_{\mathrm B}C_\lambda\le b$.
This argument does not use countable stability of box dimension.
Together with \eqref{typical:hausdorff-lower}, the inclusion
$S_\lambda\subseteq C_\lambda$, and the inequalities
\[
 \dimH\le\dimP\le\overline{\dim}_{\mathrm B},
 \qquad
 \dimH\le\underline{\dim}_{\mathrm B}
                \le\overline{\dim}_{\mathrm B},
\]
it proves \eqref{typical:four-dimensions}. Finally, if
$\lambda\equiv2\pmod3$ and $x\in\Sigma\setminus\{0\}$,
write $x=3^v u$ with $u\equiv1\pmod3$. The first nonzero digit
of $\lambda x$ is $2$, proving the remaining assertion of
Theorem~\ref{typical:main}.

\begin{corollary}[The canonical measure has smaller dimension]
\label{typical:canonical-measure}
On $S_\lambda$, start with digit $1$, choose each of the two
admissible digits with probability $1/2$ at every free step, and
take the forced digit at every other step. For almost every
$\lambda\in H$ this probability measure has full support $S_\lambda$
and is exact-dimensional of dimension $(\log_3 2)/3$, strictly
less than $\dimH S_\lambda=\log_3(4/3)$.
\end{corollary}

\begin{proof}
This is the measure in Lemma~\ref{typical:block-measure} with $q=1$.
The initial prefix is $1$, and every subsequent admissible cylinder
has positive mass, so the support is all of $S_\lambda$.
For the three translations modulo $3$, the counts $R_1$ are
$2,1,1$. Thus $H_1=(\log2)/3$, giving the stated dimension.
The strict inequality follows from $(4/3)^3>2$.
\end{proof}

\subsection{All unit parameters: a sharp gap and oscillation}

The almost-everywhere theorem leaves individual parameters
unclassified. A complementary uniform bound follows by extending
the digit coding of Abram, Bolshakov and Lagarias
\cite[Theorem~6.2]{ABL} from integer multipliers to arbitrary
$3$-adic units. The extension below is isometric and controls
cylinder counts uniformly, giving an Assouad bound.

\begin{theorem}[Isometric coding and sharp dimension gap]
\label{typical:golden-gap}
Put $\varphi=(1+\sqrt5)/2$ and $g=\log_3\varphi$.
If $\lambda\in H\setminus\{1\}$ and $k=v_3(\lambda-1)$, then
$C_\lambda$ admits an isometric embedding into
\[
 \mathcal H_k=\left\{\sum_{j\ge0}b_j3^j:
       b_j\in\{0,1\},\ b_jb_{j+k}=0\text{ for every }j\ge0\right\}.
\]
Consequently $\dim_{\mathrm A}C_\lambda\le g$. The bound is sharp:
$C_{1+3^k}=\mathcal H_k$ has Hausdorff, packing, lower and upper
box, and Assouad dimensions $g$ for every $k\ge1$.
For $\lambda=1-3^k$, all these dimensions are zero.
\end{theorem}

\begin{proof}
Write $\lambda=1+3^ku$ with $u\in\Z_3^\times$.
For an admissible input prefix $P_j=\sum_{i<j}a_i3^i$ set
\[
 r_j=\frac{\lambda P_j-[\lambda P_j]_j}{3^j},
 \qquad d_j=[r_j]_1.
\]
If $d_j=0$ both input digits are allowed; if $d_j=1$ only $0$
is allowed; if $d_j=2$ only $1$ is allowed. Define, for $r\in\Z_3$,
\[
 F(r)=\frac{r-[r]_1}{3}+\boldsymbol1_{\{[r]_1=2\}}.
\]
Every admissible step satisfies
\begin{equation}\label{typical:carry-recursion}
 r_{j+1}=F(r_j)+3^{k-1}ua_j,
 \qquad
 F(r+3^vt)=F(r)+3^{v-1}t\quad(v\ge1).
\end{equation}
At a free position $j$ and for either alternative $a\in\{0,1\}$,
put
\[
 \rho_j(a)=\big[F^{k-1}(F(r_j)+3^{k-1}ua)\big]_1.
\]
This is the residue of $r_{j+k}$ for every admissible choice of the
intervening digits. Indeed, a digit inserted strictly after $j$
contributes $3^{k-1}u$ times that digit and undergoes fewer than
$k-1$ further applications of $F$ before that residue is read.
By \eqref{typical:carry-recursion} its contribution is still a
multiple of $3$. Comparing the trajectories one insertion at a time
proves the assertion, including $k=1$ where there are no intervening
digits. The two choices at $j$ instead give
\[
 \rho_j(1)-\rho_j(0)\equiv u\not\equiv0\pmod3.
\]
Thus at most one alternative makes position $j+k$ free.

Encode a forced position by $b_j=0$. At a free position, if exactly
one alternative has $\rho_j(a)=0$, encode that alternative by $0$
and the other by $1$; if neither has zero residue, put $b_j=a_j$.
Whenever $b_j=1$, position $j+k$ is forced, so $b_{j+k}=0$.
The code belongs to $\mathcal H_k$ and depends only on the input
prefix through $j$. If two inputs first differ at $j$, their common
past makes position $j$ free, and its two choices receive different
code digits. Their first differing position is preserved exactly,
which proves isometry.

Within $L$ consecutive positions the constraint defining
$\mathcal H_k$ separates into $k$ residue classes. If class $r$
occurs $\ell_r$ times, it has $F_{\ell_r+2}$ words without adjacent
ones, where $F_0=0,F_1=1$ are the Fibonacci numbers. The same product
is an upper bound after any fixed prefix, since boundary constraints
can only remove possibilities. As $F_{\ell+2}\le\varphi^{\ell+1}$,
the number of continuations is at most $\varphi^{L+k}$. This uniform
cylinder estimate gives $\dim_{\mathrm A}\mathcal H_k\le g$.
There is no carry in $x+3^kx$ for $x\in\Sigma$, so the condition
that it belong to $\Sigma$ is exactly $a_ja_{j+k}=0$.
Hence $C_{1+3^k}=\mathcal H_k$. Its Hausdorff dimension $g$ is the
known family computation of \cite[Theorems~2.3 and~5.5]{AL2014};
combined with the Assouad upper bound, this gives all the stated
dimension equalities.

For $\lambda=1-3^k$, subtract $3^kx$ digit by digit from $x$.
If a first borrow occurs, its output digit is $2$, which is
forbidden. Thus admissibility is equivalent to
$a_j\ge a_{j-k}$ for every $j\ge k$. Along each residue class
modulo $k$ the sequence is $0^t1^\infty$, with $t\ge0$, or
$0^\infty$. This gives a countable set. More strongly, after any
fixed prefix the number of length-$L$ continuations is at most
$(L+1)^k$, which is bounded by $C_{\varepsilon,k}3^{\varepsilon L}$
for every $\varepsilon>0$. Its Assouad dimension, and hence all
the listed dimensions, are zero.
\end{proof}

\begin{corollary}[Exact punctured limits at the identity]
\label{typical:oscillation}
For $f(\lambda)=\dimH C_\lambda$, in the $3$-adic topology,
\[
 \liminf_{\substack{\lambda\to1\\\lambda\ne1}}f(\lambda)=0,
 \qquad
 \limsup_{\substack{\lambda\to1\\\lambda\ne1}}f(\lambda)
      =\log_3\varphi<\log_3 2=f(1).
\]
The same conclusions hold with packing, either box dimension,
or Assouad dimension in place of Hausdorff dimension. No unit
parameter has a dimension in $(\log_3\varphi,\log_3 2)$, and
only $\lambda=1$ has dimension $\log_3 2$.
\end{corollary}

\begin{proof}
The sequences $1-3^k$ and $1+3^k$ tend to $1$ and attain the
two limiting values, while Theorem~\ref{typical:golden-gap}
provides the uniform upper bound away from $1$ in $H$.
Parameters in $2+3\Z_3$ have intersection $\{0\}$, and
$C_1=\Sigma$ has all the stated dimensions $\log_3 2$.
\end{proof}

\begin{remark}[Scope of the parameter theorems]\label{typical:scope}
Theorem~\ref{typical:main} concerns almost every multiplier in an
uncountable parameter space. It gives no conclusion about a
prescribed integer multiplier or the sequence $4^n$. Theorem
\ref{typical:golden-gap} is uniform but is an upper bound, not a
classification of the spectrum below $\log_3\varphi$.
Neither statement decides $1\notin E_2$.
\end{remark}
